\chapter{Media Pipelines: Budget Relays Made Concrete}
\label{ch:pipelines}

\begin{quotation}
\noindent\textit{Synopsis.} This chapter applies the full framework to a
real transformation chain: audio recording, signal processing,
transcription, timestamp generation, summarization, and archiving, each
analyzed as a stage in a budget relay preserving some distinctions while
discarding others according to local priorities. The chapter provides a
concrete, constructive witness for the Composition Failure Theorem by
exhibiting a distinction lost early in the chain that becomes critical for
a later task no stage anticipated, and closes with practical
recommendations for provenance design, root-object preservation, and
re-derivable archival systems.
\end{quotation}

\section{The Full Worked Chain, Named and Numbered}

The chain audio $\to$ transcript $\to$ summary $\to$ archive is mapped
directly onto \cref{def:substrate-chain}'s notation, but rather than
leave the distinctions involved abstract, as
Appendix~\ref{app:proof}'s canonical witness does, this section fixes
concrete, named ones, so that the constructive witness promised by this
chapter's synopsis is worked with real content rather than only with
placeholders standing in for it.

\begin{construction}[Named Pipeline Instance]
\label{con:pipeline-instance}
Let $A_0$ carry four named distinctions from a recorded utterance:
lexical content $d_{\mathrm{lex}}$ (cost $3$), timestamp precision
$d_{\mathrm{time}}$ (cost $1$), sentence-level structure
$d_{\mathrm{struct}}$ (cost $2$), and emotional tone $d_{\mathrm{tone}}$
(cost $4$). The transcription stage has task $\tau_0$ requiring exactly
$\{d_{\mathrm{lex}}, d_{\mathrm{time}}\}$, with budget $\beta_0 = 4$. The
summarization stage has task $\tau_1$ requiring exactly
$\{d_{\mathrm{lex}}\}$, with budget $\beta_1 = 3$.
\end{construction}

\begin{proposition}[The Named Instance Exhibits Composition Failure]
\label{prop:named-instance-failure}
In \cref{con:pipeline-instance}, $d_{\mathrm{tone}}$ is dropped at stage
$0$ under any construction satisfying \cref{def:budget-discard}, and the
resulting archive is insufficient for any task recovering
$d_{\mathrm{tone}}$.
\end{proposition}

\begin{proof}
The required set at stage $0$, $\{d_{\mathrm{lex}}, d_{\mathrm{time}}\}$,
has cost $R_0 = 3 + 1 = 4 = \beta_0$ exactly. The non-required
distinctions have costs $C(d_{\mathrm{struct}}) = 2 > \beta_0 - R_0 = 0$
and $C(d_{\mathrm{tone}}) = 4 > 0$; by \cref{lem:exhaustion}, both are
dropped under any construction, so $D_0 = \{d_{\mathrm{struct}},
d_{\mathrm{tone}}\}$. At stage $1$, $A_1$ retains only
$\{d_{\mathrm{lex}}, d_{\mathrm{time}}\}$; the required set
$\{d_{\mathrm{lex}}\}$ has cost $R_1 = 3 = \beta_1$ exactly, and the sole
non-required distinction present, $d_{\mathrm{time}}$, has cost $1 >
\beta_1 - R_1 = 0$, so $D_1 = \{d_{\mathrm{time}}\}$ by the same lemma.
$A_2$ retains only $\{d_{\mathrm{lex}}\}$. Let $\tau^* = \tau_{d_{\mathrm{tone}}}$,
the recovery task (\cref{def:recovery-task}) for whether the speaker's
tone on a given claim was, for instance, sincere or sarcastic. Since
$d_{\mathrm{tone}} \in D_0 \subseteq D(A_0 \Rightarrow A_2)$, by
\cref{lem:recovery-sufficiency} $A_2$ is insufficient for $\tau^*$.
Neither $\tau_0$ nor $\tau_1$ is $\tau^*$: both are preservation tasks
over different, disjoint required sets not including
$d_{\mathrm{tone}}$'s recovery.
\end{proof}

This is Appendix~\ref{app:proof}'s abstract construction re-run with
distinctions a reader can recognize without translation: a transcription
system tuned for lexical accuracy and rough timing, and a summarizer
tuned for gist, jointly and irreversibly discard whether a speaker meant
something sincerely or sarcastically -- not through negligence at either
stage, but because neither stage's task ever prioritized it, exactly as
\cref{thm:composition-failure} predicts.

\section{Fan-Out, Concretized}

\cref{con:pipeline-instance} extends directly to the fan-out case of
Chapter~\ref{ch:fanout} by adding a second branch from $A_0$: an
acoustic-tone-analysis branch $Y$ with task $\tau_Y$ requiring exactly
$\{d_{\mathrm{tone}}\}$, budget $\beta_Y = 4$ (funding $d_{\mathrm{tone}}$
exactly), dropping everything else. This is
\cref{con:two-branch} instantiated with named distinctions: branch $X$
(transcription, as above) and branch $Y$ (tone analysis) are both
locally admissible, both derived from the identical $A_0$, and disagree
silently in precisely \cref{def:silent-disagreement}'s sense -- a
consumer relying on the transcript alone has no signal that tone was
ever at issue, while branch $Y$, unconsulted, would have flagged it
directly.

\begin{proposition}[Pipeline Witness]
\label{prop:pipeline-witness}
There exists an explicit audio $\to$ transcript $\to$ summary chain and a
task $\tau^*$, constructible without appeal to hypothetical or idealized
stages, such that $\tau^*$ requires a distinction discarded upstream and
no downstream stage's loss ledger records enough about that distinction
for a consumer pursuing $\tau^*$ to detect the gap.
\end{proposition}

\begin{proofsketch}
\cref{con:pipeline-instance} together with
\cref{prop:named-instance-failure} constitutes exactly this witness, with
$\tau^* = \tau_{d_{\mathrm{tone}}}$ constructed explicitly rather than
posited. This result was labeled a Theorem in an earlier draft; the
label has been corrected to Proposition, consistent with the monograph's
epistemic-labeling discipline, since it is a single explicit
instantiation of the general result already proved in
Appendix~\ref{app:proof} rather than an independent derivation.
\end{proofsketch}

\section{Practical Recommendations}

Preserving the raw audio as a root object (\cref{def:root-object}) and
treating every downstream rendering as disposable and re-derivable
directly operationalizes \cref{prop:root-recovery}: a discovered
insufficiency at the summary or transcript stage can then be repaired by
re-deriving a new rendering from the root under a revised discard
policy, rather than requiring recovery of information that no longer
exists anywhere in the chain. In terms of \cref{con:pipeline-instance},
this means that discovering $\tau_{d_{\mathrm{tone}}}$'s relevance after
the fact is recoverable, at the cost of reprocessing, if and only if the
raw audio was retained; it is permanently unrecoverable if only the
transcript or summary survived.

A budget-aware loss ledger (\cref{def:loss-ledger}) attached to each
stage should record not merely that $d_{\mathrm{struct}}$ and
$d_{\mathrm{tone}}$ were dropped at transcription, but the budget
$\beta_0 = 4$ and required set $\{d_{\mathrm{lex}}, d_{\mathrm{time}}\}$
that produced the drop, so a later archivist can distinguish a
distinction dropped because it was judged irrelevant from one dropped
purely because the transcription budget was small -- exactly the
distinction \cref{def:loss-ledger} was built to preserve. Where a
fan-out structure exists, as in Section~\ref{sec:coord-provenance}'s
sense of provenance-sharing extended across branches rather than only
along a single chain, cross-referencing each branch's ledger against the
others makes \cref{def:silent-disagreement}'s failure mode detectable
after the fact, even though, as \cref{prop:provenance-insufficient}
showed, detectability alone does not prevent the original loss.
